\documentclass[../DoAn.tex]{subfiles}
\begin{document}
{%
\hyphenpenalty=10000
\exhyphenpenalty=10000

\begin{center}
    \Large{\textbf{ABSTRACT}}\\
\end{center}
\vspace{1cm}

Remotely Operated Vehicles (ROVs) are increasingly deployed in underwater survey missions including the inspection of submerged structures, marine ecosystem monitoring, and geological data collection. In practice, all sensor data (temperature, pressure, depth, GPS positioning) and multimedia content (video, images) are recorded locally on the vehicle, and uploaded by operators only after the ROV returns to shore due to the absence of stable network connectivity in field environments. The post-dive analysis workflow currently relies entirely on manual cross-referencing: analysts must simultaneously open a video player and a CSV spreadsheet, manually correlating each data row with the corresponding video frame to determine environmental conditions at a specific moment. This approach is not only time-consuming and error-prone, but also lacks any automation capability and makes collaborative analysis within an operations team impractical. Existing solutions such as QGroundControl and BlueOS are designed primarily for real-time vehicle control and do not address the need for multi-user, web-based post-mission data management and analysis.

This thesis proposes and implements a web-based system titled \textit{``AI-Integrated ROV Operations Management''} following a multi-tier Client-Server architecture, with the digitization and automation of post-mission analysis as its primary focus. The technology stack consists of Node.js/Express.js for the server tier, MongoDB for flexible document-oriented sensor data storage, React 18 with Vite for the interactive user interface, and Redis for task queue management and session security. The most distinctive aspect of the system compared to existing solutions is the integration of three independent artificial intelligence modules: an object detection module using the YOLOv8n model running on a Python FastAPI microservice, which automatically detects and tracks objects in video frames on a per-frame basis via an asynchronous Bull Queue pipeline; a mission summarization module leveraging the Gemini 2.5 Flash API to synthesize all dive information into a natural language report; and a sensor anomaly detection module implementing the Z-Score algorithm without requiring any labeled training data to automatically flag abnormal sensor readings. The system additionally provides real-time video-to-sensor-chart synchronization and an Evidence System that allows operators to capture frame snapshots and mark video clip segments directly while reviewing footage.

The system follows a proposed deployment model using Docker Compose with four services (nginx, node-backend, yolo-service, redis) on a VPS, with the frontend distributed via Vercel CDN, the database hosted on MongoDB Atlas, and file storage on AWS S3. Functional testing results show 37 out of 37 test cases passing (100\%), covering authentication, role-based access control (RBAC), input validation, and full CRUD operations across all resources. Load testing using a custom Node.js script at 20 concurrent users demonstrated average response times below 500ms with a p95 below 1 second for standard data endpoints; the aggregation-heavy dashboard statistics endpoint achieved a p95 of approximately 1.8 seconds owing to seven parallel MongoDB aggregation pipelines. The YOLOv8n model processes a single image in under one second and handles video with adaptive frame sampling (0.2--1.0 seconds per sample depending on clip duration). All AI analysis results are pushed to the client interface via Server-Sent Events (SSE) without requiring page reloads, providing a seamless real-time experience for the operator.

The system establishes a foundation for the comprehensive digitization of ROV operations and data management workflows, replacing fragmented manual processes with a centralized, multi-user, and automated platform. In terms of future development, the system can be extended to integrate a real-time Ground Control Station (GCS) stream when network connectivity is available in the field, to fine-tune the YOLOv8 model on domain-specific ROV datasets (coral reefs, debris classification, marine species) for improved detection accuracy, or to develop a companion mobile application for on-site operator use. The microservice architecture already in place provides a natural extension point for incorporating additional specialized AI models without restructuring the existing system.

}% end hyphenpenalty group
\end{document}
